\section{Convergence of the coupling and activity histories}
\label{sec:coupled-background-convergence}
We next prove convergence of the background on which the marked
comparison depends. The coupling and the activities must be treated
together. The activity recursion contracts towards coarser scales,
whereas the coupling is fixed at the bottom scale. A finite telescoping
identity joins these two directions. In particular, the proof does not
assume that the activity space has compact bounded balls.

\subsection{The normalized variables and their paid estimates}
Fix a period. Let $\widehat\beta_m>0$ be the reference coupling at
scale $m$, and set
\[
 w_m=\widehat\beta_m^{-1/32},\qquad
 B_m=\widehat\beta_m w_m=\widehat\beta_m^{31/32},\qquad
 b_m(\lambda)=\widehat\beta_m(w_m\lambda)^2.
\]
A coupling $\beta$ is represented by
$\lambda=w_m^{-1}\sqrt{\beta/\widehat\beta_m}$. On the paid
band its normalized variable satisfies
\[
                  17/20\le w_m\lambda\le23/20.
\]
Thus, for two variables in that band,
\begin{equation}\label{eq:coupled-inverse-bound}
 B_m|\lambda-\lambda'|
 \le |b_m(\lambda)-b_m(\lambda')|.
\end{equation}
Indeed the difference on the right is
$B_m|\lambda-\lambda'|\,|w_m\lambda+w_m\lambda'|$, whose last
factor is at least $17/10>1$.

Let $\mathcal H_m$ be the completion of the seminormed space of
actual regular and forced activity arrays with the prescribed
weighted derivative and value prices. The embedding preserves
distances. In particular, equality of evaluated gases alone is
not the equivalence relation used here. Let $x_m$ denote the
embedded array. The finite-history construction gives
$\|x_m\|\le1/8$. For a terminal index $N$, write
\[
 x_m=F_{m,N}(\lambda_{m+1},x_{m+1}),\qquad
 q_m=Q_{m,N}(\lambda_{m+1},x_{m+1})\quad(m<N)
\]
for its activity transformation and the scalar correction to the
coupling. These maps are initially defined on the actual allowed
arrays; the notation does not assume an extension to every element
of the completion. Their output and comparison estimates are
\begin{align}
 \|F_{m,N}(\lambda,x)\|&\le1/8,
 &|Q_{m,N}(\lambda,x)|&\le s_m/2,\label{eq:coupled-output-bounds}\\
 \|F_{m,N}(\lambda,x)-F_{m,N}(\lambda',x')\|
 &\le c\max\{|\lambda-\lambda'|,\|x-x'\|\},\label{eq:coupled-activity-step}\\
 |Q_{m,N}(\lambda,x)-Q_{m,N}(\lambda',x')|
 &\le4s_m\max\{|\lambda-\lambda'|,\|x-x'\|\},\label{eq:coupled-cost-step}
\end{align}
where $c=8/128=1/16$ and $s_m\ge0$. The numerical schedule
also pays
\begin{equation}\label{eq:coupled-cumulative-payment}
                  \sum_{j<m}s_j\le\frac1{128}B_m.
\end{equation}
The coupling recursion is
\begin{equation}\label{eq:coupled-recursion}
 b_m(\lambda_m)-b_{m+1}(\lambda_{m+1})-\alpha_{m,N}=q_m,
\end{equation}
where $\alpha_{m,N}$ is the path correction, already convergent
at each fixed $m$ by Section~\ref{sec:background-path-convergence}.

We use the following precise continuity statement for the local
transformation: if two sequences of terminal indices tend to infinity,
their input normalized couplings have a common limit, and their
actual input arrays are Cauchy with difference tending to zero in
$\mathcal H_{m+1}$, then the differences of their output arrays and
scalar corrections tend to zero. This statement concerns arbitrary
allowed inputs, including the truncated histories constructed below.
Its analytic content is the common-variable comparison of
Section~\ref{sec:marked-kernel-convergence} without a mark, together
with the weighted head-and-tail argument of
Section~\ref{sec:local-source-stability}. Here is the precise tail implication. For a regular entry put
$b_f(i)=\sup_U|f_i^{\rm reg}(U)|$, with its actual mask, and for
a forced entry put $b_f(i)=\sup_U|f_i^{\rm forced}(U)|$.
Summing the site's paid budgets gives, separately for these two
families,
\[
 \sum_i b_f(i)\le C_m\|f\|,\qquad
 b_f(i)\le b_{f-g}(i)+b_g(i).
\]
The constants are the finite number of sites times the corresponding
positive regular or forced normalization. Fix one sufficiently late
array $f_{N_0}$ and choose a finite set outside which its $b$-sum
is small. For every later array, its tail sum is at most that fixed
tail plus $C_m\|f_N-f_{N_0}\|$. Enlarge the finite set for the
finitely many earlier arrays. This gives the uniform value-price
tails needed by the compensated comparison. It does not assert
that arbitrary activity arrays can be truncated in their full
derivative norm.

For the scalar output, use convergence of the finite projected
coupling coefficient itself: its projection is bounded by the
local derivative price, whose convergence is supplied by the
compensated comparison. This is in addition to convergence of
the output activity array. Subtracting the convergent incoming
coupling and path correction gives continuity of $Q$ by
\eqref{eq:coupled-recursion}.
The bounds \eqref{eq:coupled-output-bounds}--\eqref{eq:coupled-cost-step}
are the quantitative one-step hypotheses used in this section;
their purchase is kept separate from the convergence proof.

\subsection{First fix the normalized coupling history}
Suppose two sequences of allowed normalized coupling histories
$\lambda_{m,N}$ and $\lambda'_{m,N'}$ have the same limit for
every fixed $m$. For an integer $h\ge m$, construct a truncated
activity history by putting $x_h^{[h]}=0$ and applying the same
finite-history maps backwards down to $m$. This construction is
allowed for every prescribed coupling history in the band:
\eqref{eq:coupled-output-bounds} keeps all its arrays in the ball
of radius $1/8$. It does not impose the coupling recursion on this
auxiliary history.

Compare it with the original activity history for the same
terminal index and the same prescribed couplings. At scale $h$
the difference is at most $1/4\le1$. By
\eqref{eq:coupled-activity-step}, backwards induction gives
\begin{equation}\label{eq:coupled-truncation}
                         \|x_m-x_m^{[h]}\|\le c^{h-m}.
\end{equation}
For fixed $h$, the two truncated histories have common limits
at every scale at most $h$. At $h$ they are both zero. If this
is known at $i+1$, their arrays are Cauchy and their difference
tends to zero there; the local continuity statement and the
common normalized coupling limit give output differences tending
to zero at $i$. To get convergence, apply this comparison to
arbitrary pairs of subsequences of either truncated history.
If that history were not Cauchy, one could select indices
$u_n,v_n\ge n$ whose distances stay above some $\varepsilon>0$,
contradicting the comparison. Completeness of $\mathcal H_i$
then gives a limit, common to the two histories. This completes
the finite backwards induction. It also proves common limits
of their scalar corrections at scales below $h$.

Now choose $h-m$ so large that $2c^{h-m}<\varepsilon/2$.
Equation~\eqref{eq:coupled-truncation} and the convergence of
the two truncated histories show that the original activity
difference is eventually below $\varepsilon$. Apply the same
arbitrary-subsequence argument once more. The original activity
histories converge in $\mathcal H_m$ to the same limit. Their
scalar corrections then have common limits by local continuity
at the next scale. Thus convergence of all the normalized
couplings implies convergence of all the activities and costs.

\subsection{Identifying the possible coupling limits}
At each fixed scale the band for $\lambda_{m,N}$ is a compact
real interval. Successive subsequence selection at scales
$0,1,\ldots$ followed by the diagonal subsequence gives a
subsequence on which every normalized coupling converges.
The preceding argument attaches activity and scalar-correction
limits to it. Denote one such limiting history by
$(\lambda_m,x_m,q_m)$ and another by
$(\lambda'_m,x'_m,q'_m)$. Suppose the bottom couplings of the
two sequences have a common limit. Passing to the limit in
\eqref{eq:coupled-recursion} and subtracting gives the exact
finite identity
\begin{equation}\label{eq:coupled-telescope}
 b_m(\lambda_m)-b_m(\lambda'_m)
                  =-\sum_{j<m}(q_j-q'_j).
\end{equation}
The common limits of $\alpha_{j,N}$ cancel. There is no infinite
sum and no terminal boundary term in this identity.

The one-step inequalities also pass to these limits, despite
being initially defined only for actual input arrays. To see
this without postulating an extension of $F_{m,N}$, construct
at each terminal index a new backwards activity history with
the fixed normalized coupling script $(\lambda'_i)_i$ and zero
terminal array. Its couplings are allowed, because the band is
closed. By the previous subsection, its activities and costs
have precisely the limits $x'_m,q'_m$. Compare it to the first
history at the same terminal index using
\eqref{eq:coupled-activity-step} and
\eqref{eq:coupled-cost-step}, then pass to the limit. Writing
\[
       d_m=\max\{|\lambda_m-\lambda'_m|,\|x_m-x'_m\|\},
\]
we obtain
\begin{equation}\label{eq:coupled-limit-comparison}
 \|x_m-x'_m\|\le c d_{m+1},\qquad
                         |q_m-q'_m|\le4s_m d_{m+1}.
\end{equation}

First observe that $d_m\le1/4$ for all $m$. The activity bound
follows from $\|x_m\|,\|x'_m\|\le1/8$. For the normalized
coupling, \eqref{eq:coupled-inverse-bound},
\eqref{eq:coupled-telescope}, the bound $|q_j|,|q'_j|\le s_j/2$
and \eqref{eq:coupled-cumulative-payment} give
\[
 B_m|\lambda_m-\lambda'_m|\le\sum_{j<m}s_j
                         \le B_m/128.
\]
Thus $D=\sup_m d_m$ exists and lies in $[0,1/4]$.
Using the sharper cost comparison now gives
\[
 B_m|\lambda_m-\lambda'_m|
 \le4\sum_{j<m}s_jd_{j+1}
 \le\frac{4D}{128}B_m.
\]
Together with \eqref{eq:coupled-limit-comparison} this yields
$d_m\le(8/128)D$ for every $m$. Taking the supremum,
$0\le D\le D/16$, hence $D=0$. Both limiting normalized
coupling and activity histories are identical. The costs are
identical too, by their one-step comparison.

Every subsequence has a further subsequence with all normalized
couplings convergent, and all such limits have just been identified.
If the original normalized coupling at some scale failed to
converge to this value, a subsequence would remain outside one
fixed neighborhood, contradicting its convergent further
subsequence. Thus the entire normalized coupling sequence
converges at every scale. The first part of the proof then gives
convergence of the entire activity and correction sequences.
In particular this applies when the bottom coupling is fixed
exactly, as in the selected public construction.

\subsection{Why the cumulative payment can be imposed}
The numerical inequality \eqref{eq:coupled-cumulative-payment}
is a finite cumulative bound, not summability of a fixed sequence
of costs over all scales. We record its purchase to distinguish
the two assertions. On a stage with positive reference increment
$a$, put $B(x)=x^{31/32}$. For $x\ge a>0$, the mean value theorem
and monotonicity of $x^{-1/32}$ give
\[
 B(x+a)-B(x)\ge a\frac{31}{32}2^{-1/32}x^{-1/32}.
\]
The corresponding cost allowance is twice the activity allowance
divided by the squared fluctuation threshold. Its paid form is
\[
 s(x)=\frac{2vC_{\rm st}x^{-11/10}}{t(x)^2},\qquad
 t(x)=x^{-1/2}\exp((\log\log x)^2),
\]
where $v,C_{\rm st}>0$ are fixed at this stage. For sufficiently
large $x$, $t(x)\ge x^{-1/2}$, so
$s(x)\le2vC_{\rm st}x^{-1/10}$. Since
$-1/10<-1/32$, increasing the entry threshold gives
\[
                s(x)\le\frac1{128}\{B(x+a)-B(x)\}.
\]
There are only finitely many permitted next stages at each stage,
so their thresholds can be increased together after their block
lengths, positive increments and constants have been fixed.
The reference recursion is
$\widehat\beta_{m+1}=\widehat\beta_m+a_m$.
Summing the preceding inequalities over $j<m$ telescopes and gives
\[
 \sum_{j<m}s_j\le\frac1{128}(B_m-B_0)\le\frac1{128}B_m.
\]
This proves the cumulative estimate required in the uniqueness
argument. It preserves the chosen maps and block lengths, just
as the other entry-threshold increases do.
